\documentclass{article}
% Source: Topic_09_Teacher_Edition.pdf, PDF pages 12-24 (printed pages 443A-451B)
\usepackage{amsmath}
\usepackage{amssymb}
\usepackage{graphicx}
\usepackage{tikz}
\usepackage{pgfplots}
\pgfplotsset{compat=1.18}
\usepackage{booktabs}
\usepackage{xcolor}

\newenvironment{lesson-meta}{}{}        % lesson code, title, objective, EQ, vocab
\newenvironment{item-analysis}{}{}      % Savvas's Example × DOK table
\newenvironment{model-discuss}{}{}      % Launch opener
\newenvironment{concept-box}{}{}        % in-lesson concept reference
\newenvironment{concept-summary}{}{}    % end-of-lesson summary
\newenvironment{example}[1]{}{}         % #1 = example number
\newenvironment{tryit}[1]{}{}           % #1 = tryit number
\newenvironment{practice}[2]{}{}        % #1 = practice number, #2 = DOK
\newenvironment{te-addendum}[2]{}{}     % #1 = anchor example, #2 = type
\newcommand{\answer}[1]{}               % answer key, inside any env
\newcommand{\placeholder}[2]{}          % #1 = type, #2 = description
\newcommand{\uncertain}[2]{#1}          % #1 = best guess, #2 = alternatives
\newcommand{\te}[1]{}                   % teacher-only inline annotation

\begin{document}
% Source page 12: 443A Lesson Overview
\begin{lesson-meta}
lesson: 9-1
title: Parabolas
standards: none printed
practices: none printed
objective: I can understand and apply the geometric properties of a parabola.

essential question: What are the geometric properties of a parabola, and how do they relate to algebraic representations of a parabola?

\textbf{VOCABULARY}
\item conic section
\item directrix
\item focal length
\item focus of a parabola
\item general form of a second-degree equation
\textbf{TOPIC VOCABULARY}
\te{No topic vocabulary list is printed on these pages.}
\begin{tabular}{ll}
Lesson Code & 9-1 \\
Title & Parabolas \\
Standards & none printed \\
I CAN Objective & I can understand and apply the geometric properties of a parabola. \\
Essential Question & What are the geometric properties of a parabola, and how do they relate to algebraic representations of a parabola? \\
Lesson Vocabulary & conic section; directrix; focal length; focus of a parabola; general form of a second-degree equation \\
Topic Vocabulary & none printed \\
\end{tabular}
\end{lesson-meta}
\begin{te-addendum}{0}{GeneralizeETP}
\textbf{Lesson Overview: FOCUS}
Objective. Students will be able to:
\item Derive the equation of a parabola.
\item Relate a parabola's focal length to its equation.
\item Rewrite an expression by completing the square and then use it to find the focus and directrix of a parabola.
Essential Understanding. A parabola is the set of all points on a plane that are equidistant from the focus and the directrix. This definition is used to derive the equation of a parabola and apply the geometric properties of a parabola.
\textbf{COHERENCE}
Previously in course, students:
\item Understood a parabola as the shape created by the graph of a quadratic function.
In this lesson, students:
\item Understand that a parabola is a type of conic section.
\item Derive and write the equation of a parabola.
Later in this topic, students will:
\item Derive and write the equations of other conic sections including circles, ellipses, and hyperbolas.
\textbf{RIGOR}
This lesson emphasizes a blend of conceptual understanding and application.
\item Students understand that the definition of a parabola can be used to derive the equation of a parabola because all points on the parabola are equidistant from the focus and directrix.
\item Students apply the equation of a parabola to model real-world situations, such as identifying the location for placing a bracket on a dish with a parabolic shape.
\textbf{Mathematics Overview}
Students begin their investigation of conic sections with parabolas. They derive the equation of a parabola using the focus and directrix or the focal length of the parabola. Students also rewrite equations in the vertex form, completing the square where necessary, to determine the equation of a parabola.
\textbf{Applying Math Practices}
Look for and Make Use of Structure
Students use the structure of the equation of a parabola to identify key features of the graph of the parabola.
Look For and Express Regularity in Repeated Reasoning
Students notice relationships between the way a plane intersects a double right cone and the type, size, and shape of the conic section that is created.
\te{Program Resources. SavvasRealize.com.}
\end{te-addendum}
\begin{te-addendum}{0}{ELLAddendum}
\textbf{Vocabulary Builder}
Glossary is available at SavvasRealize.com.
NEW VOCABULARY
\begin{tabular}{ll}
English & Spanish \\
conic section & sección cónica \\
directrix & directriz \\
focal length & distancia focal \\
focus of a parabola & foco de una parábola \\
general form of a second-degree equation & forma general de una ecuación de segundo grada \\
\end{tabular}
\te{Printed Spanish ends in grada; likely grado.}
VOCABULARY ACTIVITY
Ask students to draw a parabola and identify the key features of the graph, such as the vertex and axis of symmetry. Introduce students to the definition of parabola used in this lesson, using a diagram to show both the focus and the directrix. Explain to students that they will use this definition to derive the equation of a parabola. Have students answer the multiple-choice question.
Every point on a parabola is equidistant from the focus and what line?
A. x-axis
B. y-axis
C. axis of symmetry
D. directrix
\answer{See back of book.}
\te{No answer marking is printed in this Vocabulary Builder preview.}
\end{te-addendum}
\begin{te-addendum}{0}{LearnTogether}
\textbf{Student Companion}
Students can do their work for the lesson in their Student Companion or on Savvas Realize.
% FIG 12 1000 787 1152 967
\placeholder{illustration}{Student Companion cover, colored glowing loops on dark background; enVision Algebra 2; SAVVAS.}
\end{te-addendum}
% Source page 13: 443B Explore
\begin{te-addendum}{0}{LearnTogether}
\textbf{EXPLORE \& REASON}
INSTRUCTIONAL FOCUS Students explore their perception of what is created when a flat plane intersects a geometrical 3-dimensional figure. This investigation is used to introduce the concept of defining conic sections using the equation of a parabola.
STUDENT COMPANION Students can complete the Explore \& Reason activity in their Student Companion.
\te{Explore \& Reason is available at SavvasRealize.com. Activity.}
\end{te-addendum}
\begin{te-addendum}{0}{GeneralizeETP}
\textbf{Before: WHOLE CLASS}
Implement Tasks that Promote Reasoning and Problem Solving ETP
Q: What do you see when you look at each figure from the side?
\answer{On the square pyramid you see triangles, and on the cylinder you see a rectangle.}
\textbf{During: SMALL GROUP}
Support Productive Struggle in Learning Mathematics ETP
Q: What do you notice about the created shapes as you move the planes through the figures?
\answer{The shapes change in size.}
For Early Finishers
Q: Consider a sphere. What can you say about the circles that are created when you intersect the sphere with a vertical plane or a horizontal plane? Use a drawing to explain.
\answer{Check students' work. For any intersection, there is a center with infinitely many points equidistant from it that create a circle.}
\textbf{After: WHOLE CLASS}
Facilitate Meaningful Mathematical Discourse ETP
Q: What do you think the shapes created in each cross section have in common?
\answer{They are all similar. They have the same shape with different side lengths, and the ratios of corresponding side lengths, are equal.}
\te{Printed answer says all cross sections are similar; likely intended for parallel sections of a given type, since arbitrary sections can have different shapes.}
\end{te-addendum}
\begin{model-discuss}
\textbf{EXPLORE \& REASON}
Consider the square pyramid and the cylinder shown.
% FIG 13 908 262 1053 350
\placeholder{diagram}{Lavender square pyramid at left with solid visible edges and dashed hidden base edges and rear lateral edge; vertical lavender circular cylinder at right with elliptical top and bottom, hidden rear bottom semicircle dashed. No labels or measurements.}
\te{Digital version: Go back. You can use the pen tool to sketch the shapes asked for in parts A and B. Enter your answer. Student Edition. SAMPLE STUDENT WORK.}
A. What shapes can be created by intersecting a vertical plane with each figure? What shapes can be created by intersecting a horizontal plane with each figure?
\answer{A. vertical plane: pyramid: trapezoid (or triangle if plane intersects the top point); cylinder: rectangle. horizontal plane: pyramid: square; cylinder: circle.}
B. Make Sense and Persevere What shapes can be created by intersecting a diagonal plane with each figure?
\answer{B. pyramid: triangle; cylinder: oval}
\end{model-discuss}
\begin{te-addendum}{0}{HabitsOfMind}
Use with EXPLORE \& REASON
Look for Relationships What are other three-dimensional figures that have a cross section that is the same as either the cylinder or the pyramid? Give two examples.
\answer{A cone and a sphere both have circular cross sections; a cone also has a triangular cross section.}
\end{te-addendum}
% Source page 14: 443 Understand & Apply
\begin{te-addendum}{0}{LearnTogether}
\textbf{INTRODUCE THE ESSENTIAL QUESTION}
Establish Mathematics Goals to Focus Learning ETP
Students will apply geometric properties of parabolas to derive the equations of specific parabolas. They will investigate the application of parabolas that can be used to model real-world situations.
\te{Examples are available at SavvasRealize.com. Activity.}
\end{te-addendum}
\begin{te-addendum}{0}{PurposefulQuestions}
\textbf{CONCEPT Conic Sections}
Q: Describe what is created by the intersection of a plane and a cone.
\answer{The intersection creates a conic section that can be a circle, an ellipse, a parabola, or a hyperbola, depending on the orientation of the plane.}
Q: How can you relate the cross sections to shapes created on a coordinate plane?
\answer{They are the same. Each cross section can be graphed on a coordinate plane.}
Q: If two-dimensional figures are created on a coordinate plane by an infinite number of points that look like lines or curves, what can you say about 3-D figures?
\answer{Three-dimensional figures are created by an infinite number of planes of 2-D cross sections.}
\end{te-addendum}
\begin{concept-box}
\textbf{CONCEPT Conic Sections}
A conic section is a curve formed by the intersection of a plane and two right cones.
\begin{tabular}{llll}
Parabola & Circle & Ellipse & Hyperbola \\
The intersecting plane cuts through one side of one cone. & The intersecting plane is perpendicular to the axis. & The intersecting plane cuts through both sides of a cone. & The intersecting plane cuts through both cones. \\
\end{tabular}
% FIG 14 842 325 916 441
\placeholder{diagram}{Parabola: pale-green double right cone with tips joined and dashed vertical axis. Blue slanted plane parallel to a cone generator cuts one side of lower cone, pale-green parabolic intersection. Hidden sections dashed.}
% FIG 14 924 326 993 433
\placeholder{diagram}{Circle: matching pale-green double right cone, dashed vertical axis; blue horizontal plane perpendicular to axis cuts lower cone in a pale-green circle shown in perspective. Hidden sections dashed.}
% FIG 14 1004 326 1074 435
\placeholder{diagram}{Ellipse: matching pale-green double right cone, dashed vertical axis; tilted blue plane crosses both sides of lower cone forming pale-green ellipse. Hidden sections dashed.}
% FIG 14 1088 326 1142 465
\placeholder{diagram}{Hyperbola: matching pale-green double right cone, dashed vertical axis; blue vertical plane offset to right of axis intersects both cones. Pale-green separate curved sections on upper and lower cones.}
Conic sections graphed in a coordinate plane can be represented by a second-degree equation in two variables.
The general form of a second-degree equation is
(Ax^2+Bxy+Cy^2+Dx+Ey+F=0).
\end{concept-box}
\begin{te-addendum}{1}{PurposefulQuestions}
\textbf{Additional Examples. ADDITIONAL EXAMPLE 1}
\te{Additional Examples are available at SavvasRealize.com. Go back. ESSENTIAL QUESTION. SOLUTION.}
A portable solar dish has a parabolic reflector to catch sun light from all directions throughout the day. The instructions suggest a distance of 12.25 inches from the hanging-hook to the vertex. The vertex is located at (0,12.25). Determine the equation of the directrix and the equation of the parabola.
directrix: (y=0)
parabola:
(\sqrt{(x-0)^2+(y-24.5)^2}=\sqrt{(x-x)^2+(y-0)^2})
(x^2+(y-24.5)^2=y^2)
(x^2+y^2-49y+600.25=y^2)
(x^2+600.25=49y)
(\frac1{49}x^2+12.25=y)
% FIG 14 398 940 531 1139
\placeholder{graph}{Gray grid spacing2; black double-arrowed x,y axes, O origin; window x[-16,16],y[-16,32], x labels -12,-8,-4,4,8,12 and y -12,-8,-4,4,8,12,16,20,24,28. Orange parabola segment y=x^2/49+12.25 with closed endpoints approximately (-13,16),(13,16), filled vertex (0,12.25). Filled orange focus (0,24.5), dashed black segments to endpoints, orange label Solar directrix inside dish. Orange origin (0,0). Labels Focus (0,24.5), Vertex (0,12.25).}
\te{Printed Solar directrix label lies inside the dish above the vertex; likely it labels the reflector rather than the directrix y=0.}
\answer{directrix: (y=0); parabola: (\frac1{49}x^2+12.25=y).}
Example 1 Students can practice writing the equation of a parabola given its graph with this additional example.
Q: What values do you need to identify on the graph in order to write the equation?
\answer{the coordinates of the focus and the coordinates of a point on the directrix}
\end{te-addendum}
\begin{te-addendum}{3}{PurposefulQuestions}
\textbf{ADDITIONAL EXAMPLE 3}
\te{Go back. ESSENTIAL QUESTION. SOLUTION.}
Two parabolas open in opposite directions and are symmetric across the x-axis. Both have a vertex (0,0). Parabola A opens upward and parabola B opens downward. The focus of Parabola B is at (0,-19.25). Write an equation for each parabola.
\begin{tabular}{ll}
Parabola A & Parabola B \\
Focal length: (c=19.25) & Focal length: (c=-19.25) \\
(y=\frac1{4(19.25)}x^2) & (y=\frac1{4(-19.25)}x^2) \\
(y=\frac1{77}x^2) & (y=-\frac1{77}x^2) \\
\end{tabular}
\answer{A. (y=\frac1{77}x^2); B. (y=-\frac1{77}x^2).}
Example 3 Students write equations of multiple parabolas graphed on the same coordinate plane given a verbal description.
Q: Which location determines the sign of the focal length c?
\answer{When the parabola has a vertical line of symmetry, the value of y determines the sign of c, and when the parabola has a horizontal line of symmetry, the value of x determines the sign of c.}
\end{te-addendum}
% Source page 15: 444 Write an Equation
\begin{te-addendum}{1}{PurposefulQuestions}
\textbf{Write an Equation of a Parabola}
Pose Purposeful Questions ETP
Q: How can you define point P in the equation for the parabola?
\answer{Point P is any coordinate point on the parabola and is represented in the equation as x and y.}
Q: Why do you use the variables x and y for P instead of its actual coordinates when writing the equation?
\answer{Point P represents an infinite number of points, so you must leave them as x and y in order to write the equation for the entire parabola.}
Q: What indicates that the parabola opens vertically or horizontally?
\answer{If the directrix is horizontal, the parabola opens vertically. If the directrix is vertical, the parabola opens horizontally.}
\te{Example 2 is Powered by Desmos. Activities are available at SavvasRealize.com. Printed Example 2 in the banner above Example 1; likely Example 1.}
\end{te-addendum}
\begin{example}{1}
\textbf{Write an Equation of a Parabola}
A. What is an equation for the parabola with focus (0,2) and directrix (y=-2)? Use the definition of a parabola to write and simplify the equation.
A parabola is the set of all points in a plane equidistant from a given point called the focus of a parabola and a given line called the directrix.
% FIG 15 814 292 955 388
\placeholder{graph}{Black double-arrowed x,y axes, O origin, no grid; unit ticks, x labels -6,-4,-2,2,4,6, y2,4, window roughly x[-7,7],y[-3,6]. Red upward parabola y=x^2/8, vertex O, arrowheads at upper ends. Blue focus F(0,2), black P(x,y) near (5,3), blue Q(x,-2) directly below P on solid blue double-arrowed directrix y=-2. Red dashed equal-distance segments F to P and P to Q labeled d, red right-angle marker at Q.}
\te{Callout: For a point P(x,y) on the parabola, the closest point Q on the directrix is (x,-2).}
By the definition of a parabola, the distance from P to F must equal the distance from P to Q.
(PF=PQ)
(\sqrt{(x-0)^2+(y-2)^2}=\sqrt{(x-x)^2+(y+2)^2})
(x^2+(y-2)^2=(y+2)^2)
(x^2+y^2-4y+4=y^2+4y+4)
(x^2=8y)
(\frac18x^2=y)
The equation for a parabola with focus (0,2) and directrix (y=-2) is (y=\frac18x^2).
\te{LOOK FOR RELATIONSHIPS What are the vertex and axis of symmetry for the parabola with equation (y=\frac18x^2)?}
\answer{A. (y=\frac18x^2)}
B. What is an equation for the parabola with focus (4,0) and directrix (x=-4)?
By sketching the parabola based on the given information, you can see that this parabola opens to the right.
% FIG 15 1014 547 1113 648
\placeholder{graph}{Unit gray grid window x,y[-5,5], black double-arrowed x,y axes, O origin, labels -4,-2,2,4. Red right-opening parabola x=y^2/16, vertex O, arrowheads at top/bottom. Solid blue double-arrowed directrix x=-4, blue focus F(4,0), black P(x,y) near (0.6,-2.5), blue D(-4,y) directly left on directrix. Red dashed segments D to P and P to F.}
Use the Distance Formula to set (PF=PD).
(PF=PD)
(\sqrt{(x-4)^2+(y-0)^2}=\sqrt{(x+4)^2+(y-y)^2})
((x-4)^2+y^2=(x+4)^2)
(x^2-8x+16+y^2=x^2+8x+16)
(y^2=16x)
(x=\frac1{16}y^2)
The equation for a parabola with focus (4,0) and directrix (x=-4) is (x=\frac1{16}y^2).
\answer{B. (x=\frac1{16}y^2)}
\end{example}
\begin{tryit}{1}
What is an equation for the parabola with the given focus and directrix?
a. focus (0,-3) and directrix (y=3)
b. focus (-2,0) and directrix (x=2)
\answer{a. (y=-\frac1{12}x^2) b. (x=-\frac18y^2)}
\end{tryit}
\begin{te-addendum}{1}{ElicitEvidence}
Elicit and Use Evidence of Student Thinking ETP
Q: When substituting in the representation of the directrix, why do you use a specific x- or y-value and a variable?
\answer{A directrix is a horizontal or vertical line, and the variable represents the fact that it can be all real numbers for that variable.}
\end{te-addendum}
\begin{te-addendum}{1}{CommonError}
Try It! 1a Some students may substitute the -3 and the 3 in for the wrong variables. Therefore they may need a guide that will help them stay organized when setting up the equation. Have them draw a visual that shows that the equation is two distances set equal to each other; the distance from P to the focus and the distance from P to the directrix.
Define how to set up the distances:
The distance between any point P, (x,y), on the parabola and the focus (0,-3) is equal to the distance between P and the directrix (x,3). Use the distance formula to show equality between the distances to create the equation.
\end{te-addendum}
% Source page 16: 445 Derive the Equation
\begin{te-addendum}{2}{GeneralizeETP}
\textbf{Derive the Equation for a Parabola}
Build Procedural Fluency From Conceptual Understanding ETP
Q: How does a parabola change as the distance between the focus and the vertex of the parabola decreases?
\answer{The parabola is stretched vertically.}
Q: For a vertical parabola with vertex at the origin, how do you identify the focus and the directrix to use when setting up the distance formulas?
\answer{Use the graph of a vertical parabola; the focus is located at (0,c) on the y-axis, and the point (x,-c) is located on the directrix.}
Q: Why does the value of c affect the width of the parabola?
\answer{Since c is a factor of the denominator of the leading coefficient, as its value moves farther from zero, the absolute value of the fraction decreases; and as its value moves toward zero, the absolute value of the fraction increases. This causes vertical compression and stretch.}
\te{Example 2 is Powered by Desmos. Activities are available at SavvasRealize.com.}
\end{te-addendum}
\begin{example}{2}
\textbf{Derive the Equation for a Parabola}
\te{CONCEPTUAL UNDERSTANDING. Powered by Desmos.}
A. What is the general equation for a parabola with a vertex at the origin and focus (0,c)?
The focal length of a parabola (|c|) is the distance between the focus and the vertex, measured along the axis of symmetry. For parabolas that open upward or to the right, (c>0). For parabolas that open downward or to the left, (c<0).
Use the focal length to find the directrix.
% FIG 16 964 329 1154 436
\placeholder{graph}{Unlabeled black double-arrowed axes, no grid/ticks. Red upward parabola with arrowheads and filled vertex (0,0); blue focus F(0,c) on positive vertical axis, black P(x,y) on right branch; blue Q(x,-c) directly below P on solid blue double-arrowed horizontal directrix below vertex. Black point where directrix crosses vertical axis. Red c labels on vertical distances above/below vertex; red dashed FP and PQ segments.}
\te{Callout: The distance from the vertex to the directrix is equal to the focal length |c|, so the directrix is y=-c.}
Based on the definition of a parabola, P(x,y) must be equidistant from the focus and the directrix. Use the Distance Formula to write an equation.
(PF=PQ)
(\sqrt{(x-0)^2+(y-c)^2}=\sqrt{(x-x)^2+(y+c)^2})
(x^2+(y-c)^2=(y+c)^2)
(x^2+y^2-2cy+c^2=y^2+2cy+c^2)
(x^2=4cy)
(\frac1{4c}x^2=y)
The equation for a vertical parabola with vertex at the origin with focal length (|c|) is (y=\frac1{4c}x^2).
The equation for a horizontal parabola with vertex at the origin with focal length (|c|) is (x=\frac1{4c}y^2).
\te{COMMUNICATE PRECISELY How could you write an equation for the directrix in terms of c?}
\answer{A. (y=\frac1{4c}x^2)}
B. How does the focal length affect the shape of a parabola?
% FIG 16 869 635 987 695
\placeholder{graph}{Black double-arrowed x,y axes, O origin, no grid; horizontal ticks every2, labels -8,-4,4,8, vertical labeled4. Three upward parabolas with vertex O and arrowheads: black narrow y=x^2/8 labeled c=2, blue middle y=x^2/16 labeled c=4, red wide y=x^2/24 labeled c=6. Black focus P(0,2), blue Q(0,4), red R(0,6); visible window about x[-10,10],y[-3,9].}
As the focal length (|c|) increases from 2 to 4 to 6, the parabola equations change from (y=\frac18x^2) to (y=\frac1{16}x^2) to (y=\frac1{24}x^2). Therefore, the parabolas are more compressed vertically as (|c|) increases.
\answer{B. The parabolas are more compressed vertically as (|c|) increases.}
\end{example}
\begin{tryit}{2}
A parabola has a focus of (4,0) and a directrix at (x=-4).
a. What is the vertex of the parabola?
b. Is the equation in the form (y=ax^2) or (x=ay^2)?
c. What is the focal length?
d. What is the equation of the parabola?
\answer{a. (0,0) b. (x=ay^2) c. 4 d. (x=\frac1{16}y^2)}
\end{tryit}
\begin{te-addendum}{2}{ElicitEvidence}
Elicit and Use Evidence of Student Thinking ETP
Q: What strategy can you use to get started?
\answer{You can use the given features to sketch the graph of the parabola on a coordinate plane.}
\end{te-addendum}
\begin{te-addendum}{2}{RtIExtend}
\textbf{Extend Student Thinking}
USE WITH EXAMPLE 2 Have students explore working with irrational numbers as focal lengths.
\item A parabola has a focus of (0,2\sqrt5) and a directrix at (y=-2\sqrt5).
Q: What is the vertex of the parabola?
\answer{(0,0)}
Q: Is this a graph of a vertical or horizontal parabola?
\answer{Vertical, the focus is on the y-axis, so the parabola is vertical.}
Q: Define the focal length. Give the exact value and an approximate value.
\answer{(2\sqrt5\approx4.47)}
Q: What is the equation of the parabola?
\answer{(y=\frac{\sqrt5}{40}x^2)}
Q: How can you sketch features of a parabola based on an irrational focal length?
\answer{Irrational numbers can be placed on a coordinate plane by estimation.}
Q: After substituting the value of c, what must be done in order to represent the parabola as an equation?
\answer{Because c is an irrational number, you have to rationalize the denominator.}
\end{te-addendum}
% Source page 17: 446 Write an Equation
\begin{te-addendum}{3}{PurposefulQuestions}
\textbf{Write an Equation of a Parabola}
Pose Purposeful Questions ETP
Q: The equation for any parabola is (y=\frac1{4c}x^2). If you are given the value of c for a specific parabola, how do you write its equation?
\answer{You can substitute the given value for c into the equation and simplify.}
\te{Printed any parabola; likely a vertical parabola with vertex at the origin is intended.}
Q: How do you know when to use (y=\frac1{4c}x^2) versus (x=\frac1{4c}y^2)?
\answer{If the parabola has a vertical axis of symmetry, use (y=\frac1{4c}x^2). If the parabola has a horizontal axis of symmetry, use (x=\frac1{4c}y^2).}
\te{Examples are available at SavvasRealize.com.}
\end{te-addendum}
\begin{example}{3}
\textbf{Write an Equation of a Parabola}
What is the equation of each parabola?
A. A parabola with vertex (0,0) and focus (0,-0.5)
Start by sketching the parabola based on the information you are given.
% FIG 17 817 291 935 371
\placeholder{graph}{Gray quarter-unit grid, window x[-1.5,1.5],y[-1.5,0.5], black double-arrowed x,y axes, O origin, x labels -1,1, y labels -0.5,-1. Red downward parabola y=-x^2/2, vertex O, arrowheads at lower left/right. Red filled focus (0,-0.5), labeled focus.}
\te{Callout: It is helpful to sketch a graph to visualize the parabola as you are writing an equation. STUDY TIP Be sure to determine if the parabola has a vertical or horizontal axis of symmetry. Note that the focus is always inside the parabola.}
By looking at the graph, you see that the parabola has a vertical axis of symmetry, so its equation has the form (y=\frac1{4c}x^2).
The coordinates of the focus are (0,-\frac12), so (c=-\frac12). Substitute into the equation (y=\frac1{4c}x^2).
(y=\frac1{4(-\frac12)}x^2=-\frac12x^2)
This parabola has the equation (y=-\frac12x^2).
\answer{A. (y=-\frac12x^2)}
B. A parabola with vertex (0,0) and directrix (x=5)
Start by sketching the parabola based on the given information.
% FIG 17 813 525 913 624
\placeholder{graph}{Unit gray grid window x[-4,6],y[-5,5], black double-arrowed x,y axes, O origin, x labels-2,2,4,y-4,-2,2,4. Red left-opening parabola x=-y^2/20, vertex O, arrowheads at top/bottom. Solid blue double-arrowed vertical directrix x=5.}
To write an equation for the parabola, first determine the focus. The vertex is equidistant between the directrix and the focus, so if the vertex is at the origin and the directrix is at (x=5), the focus must be at (-5,0). Then (c=-5).
The parabola has a horizontal axis of symmetry, so use the equation (x=\frac1{4c}y^2), where (c=-5).
(x=\frac1{4c}y^2=\frac1{4(-5)}y^2=-\frac1{20}y^2)
The equation for this parabola is (x=-\frac1{20}y^2).
\answer{B. (x=-\frac1{20}y^2)}
\end{example}
\begin{tryit}{3}
a. What is the equation of a parabola with focus (0,6) and directrix (y=-6)? How do you know that the vertex is at (0,0)?
b. What is the equation of a parabola with focus (\frac15,0) and vertex (0,0)?
\answer{a. (y=\frac1{24}x^2); The vertex is midway between the focus and the directrix, so the vertex is (0,0). b. (x=\frac54y^2)}
\end{tryit}
\begin{te-addendum}{3}{ElicitEvidence}
Elicit and Use Evidence of Student Thinking ETP
Q: By analyzing the given information, describe the graphs in Try It 3a and 3b and compare the values of a in the equations.
\answer{The graph for Try It 3a has a vertical line of symmetry and opens upward. The graph for Try It 3b has a horizontal line of symmetry and opens to the right. For Try It 3a, the value of a is (\frac1{24}). For Try It 3b, the value of a is (\frac54). Therefore, the parabola of 3a is much wider.}
\end{te-addendum}
\begin{te-addendum}{3}{HabitsOfMind}
Use with EXAMPLES 1, 2 \& 3
Communicate Precisely What is the formal geometric definition of a parabola?
\answer{A parabola is the set of all points equidistant from a given point (focus) and a line (directrix).}
\end{te-addendum}
% Source page 18: 447 Model a Real-World Situation
\begin{te-addendum}{4}{GeneralizeETP}
\textbf{Use a Parabola to Model a Real-World Situation}
Use and Connect Mathematical Representations ETP
Q: How can you relate the given equation to the equation for any parabola in terms of c?
\answer{Since the given equation is (y=\frac1{108}x^2) and the equation for any parabola is (y=\frac1{4c}x^2), you can use the Transitive Property to set the two equations equal to each other: (\frac1{108}x^2=\frac1{4c}x^2).}
Q: Since the equations are exactly the same except for the denominator, how can you solve for c?
\answer{You can set the denominators equal to each other (4c=108) and then solve for c by using the Division Property of Equality.}
\te{Examples are available at SavvasRealize.com.}
\end{te-addendum}
\begin{example}{4}
\textbf{Use a Parabola to Model a Real-World Situation}
\te{APPLICATION}
One type of solar cooker is a reflective parabolic dish. Its cross section is a parabola modeled by the equation (y=\frac1{108}x^2), with distances measured in inches. The bracket that holds the pan for the cooker needs to be placed at the focus. At what coordinates should the bracket be placed?
% FIG 18 838 296 1061 466
\placeholder{photo}{Reflective silver parabolic solar dish on yellow stand outdoors against mountains and stone wall. Cooking pan above dish center on bracket. Red point at central dish vertex with white callout Vertex:(0,0); white callout Focus:(0,c) points to pan/bracket.}
Formulate. Consider a cross-section of the parabolic dish.
% FIG 18 836 485 948 558
\placeholder{graph}{Unlabeled black double-arrowed perpendicular axes, no grid/ticks. Red upward parabola with arrowheads, red filled vertex labeled (0,0); green focus above origin on vertical axis.}
The distance from the focus to the vertex is the focal length, c.
The given equation is in the form (y=\frac1{4c}x^2), so (\frac1{108}=\frac1{4c}).
Compute. Solve for c.
(\frac1{108}=\frac1{4c})
(4c=108)
(c=27)
Interpret. In this example, the focal length is 27. The vertex is at (0,0), so the focus is at (0,27). The bracket should be placed 27 in. from the vertex.
\answer{(0,27); 27 in. from the vertex.}
\end{example}
\begin{tryit}{4}
If the reflecting dish in Example 4 has a depth of 30 in., what is the diameter of the solar cooker?
\answer{about 113.8 in.}
\end{tryit}
\begin{te-addendum}{4}{ElicitEvidence}
Elicit and Use Evidence of Student Thinking ETP
Q: If the vertex is (0,0), what do you know about the left side and the right side of the parabolic dish?
\answer{The left side and the right side of the dish are exactly the same length, and the ends are on opposite points.}
\end{te-addendum}
\begin{te-addendum}{4}{HabitsOfMind}
Use with EXAMPLE 4
Reason In Example 4, why is it helpful to locate the vertex of the parabola at the origin?
\answer{Locating the vertex at the origin produces a much simpler equation with which to find the focus.}
\end{te-addendum}
\begin{te-addendum}{4}{ELLAddendum}
\textbf{English Language Learners} (Use with EXAMPLE 4)
WRITING BEGINNING Provide students with some sentences from the example with missing words (dish, bracket, pan). Have students write sentences in their journals choosing the correct words to complete each sentence. (For example, "You cook the bacon in a pan." "You serve the meal on a dish." "You use a bracket to help build the bookcase.")
Q: What is a bracket used for?
\answer{A bracket connects two pieces of a structure together.}
Q: Is dish used in the example in the most common way the word is used?
\answer{No, a dish is usually a plate you serve food on, but in the example it describes a solar cooker in the shape of a dish (concave).}
READING DEVELOPING Have students read through the problem and underline or highlight all of the important information for solving the problem.
Q: Is it important to know what a solar cooker is? Explain.
\answer{No, it is important to know only the shape of the object.}
Q: Would you be able to solve the problem if the vertex location were not given? Explain.
\answer{No, if you did not know the location of the vertex, you could not find the exact location of the focus.}
SPEAKING EXPANDING Have students identify the adjectives used to describe the solar cooker dish, reflective and parabolic. Challenge students to identify the noun or verb within each adjective (reflect and parabola) and discuss how each word was changed into an adjective. Then have students brainstorm other adjectives that are formed in the same way (e.g., algebra/algebraic, hyperbola/hyperbolic, select/selective).
Q: What shape is an object that is described as parabolic?
\answer{a parabola}
Q: What are some ways that other shapes or math terms are changed for use as adjectives?
\answer{curve/curved, circle/circular, etc.}
\end{te-addendum}

% Source page 19: 448, Example 5
\begin{te-addendum}{5}{PurposefulQuestions}
\textbf{Pose Purposeful Questions: Complete the Square to Find the Focus and Directrix}
Q: How do you know if you need to solve for (x) or (y)?
\answer{You first determine what the quadratic expression is in the original equation.}
Q: Why is it important to complete the square first in order to find the focus and directrix?
\answer{By completing the square, you are writing the equation of the parabola in vertex form, which makes it easier to identify values to find the focus and directrix.}
Q: After you write the equation in vertex form, how do you find the value of (c)?
\answer{You already know that (a=\frac{1}{4c}). Use the vertex form of the equation to find the value of (a) and substitute this value into the equation and solve for (c).}
Q: How can the value of (c) define the focus and the directrix?
\answer{The absolute value of (c) is the distance between the vertex and the focus. The directrix is the same distance from the vertex as the focus, but in the opposite direction.}
\end{te-addendum}
\begin{example}{5}
\textbf{Complete the Square to Find the Focus and Directrix}
The equation for a parabola is given by (0=y^2-8y+x+14). Complete the square to identify the vertex, focus, and directrix.
The vertex form of an equation for a vertical parabola with vertex ((h,k)) and focal length (|c|) is ((y-k)=\frac{1}{4c}(x-h)^2). The vertex form for a horizontal parabola with vertex ((h,k)) is ((x-h)=\frac{1}{4c}(y-k)^2).
To find the vertex, focus, and directrix of the equation (0=y^2-8y+x+14), rewrite the equation by completing the square.
(0=y^2-8y+x+14)
(-x-14=y^2-8y)
(-x-14+16=y^2-8y+16)
(2-x=(y-4)^2)
(x-2=-1(y-4)^2)
\te{Callout: The equation is a horizontal parabola, in the form (x-h=\frac{1}{4c}(y-k)^2).}
The vertex of the parabola is at ((2,4)).
Since (\frac{1}{4c}=-1), solve the equation for (c) to find the focus.
(\frac{1}{4c}=-1)
(c=-\frac{1}{4})
The focus is (\frac{1}{4}) unit left of the vertex, so the focus is ((1\frac{3}{4},4)).
The directrix is also (\frac{1}{4}) unit right of the vertex, so the directrix is (x=2\frac{1}{4}).
% FIG 19 825 525 978 619
\placeholder{graph}{Black double-arrow axes x and y, O at origin; x quarter-unit ticks labeled 0.5, 1, 1.5, 2, 2.5 and y unit ticks labeled 2, 4, 6, no grid. Approximate window x [-0.4,2.7], y [-2,7]. Red left-opening parabola x=2-(y-4)^2 with vertex (2,4), arrowheads at left ends, filled red focus labeled (1 3/4,4). Blue vertical double-arrow directrix labeled x=2 1/4.}
\te{COMMON ERROR You may think that the y-coordinate of the vertex is (-4). However, remember that the term has the form ((y-k)), so in this case, (k=4), not (-4).}
\answer{Vertex ((2,4)); focus ((1\frac{3}{4},4)); directrix (x=2\frac{1}{4}).}
\end{example}
\begin{tryit}{5}
Use the vertex form to identify the vertex, focus, and directrix of the parabola with each equation.
a. (0=-y+2x^2-4x+6)
b. (0=x^2+6x+4y+5)
\answer{a. vertex: ((1,4)); focus: ((1,\frac{33}{8})); directrix: (y=\frac{31}{8}); b. vertex: ((-3,1)); focus: ((-3,0)); directrix: (y=2)}
\end{tryit}
\begin{te-addendum}{5}{HabitsOfMind}
\textbf{HABITS OF MIND: Generalize}
Q: How do you find the focus of a parabola written in vertex form?
\answer{In the equation (y=a(x-h)^2+k), (a=\frac{1}{4c}), so (c=\frac{1}{4a}), and thus the focus is ((h,k+c)).}
\end{te-addendum}
\begin{te-addendum}{5}{RtISupport}
\textbf{Support Struggling Students}
USE WITH EXAMPLE 5 Some students may struggle with setting up an equation of a parabola before completing the square.
In the following equations of a parabola, identify the quadratic expression and explain how to complete the square.
1. (0=x^2+6x-y+12)
\answer{The quadratic expression is (x^2+6x); isolate the quadratic expression, complete the square, and solve for (y).}
2. (0=-x+18+y^2-4y)
\answer{The quadratic expression is (y^2-4y); isolate the quadratic expression, complete the square, and solve for (x).}
Q: How do you recognize the quadratic expression in order to determine which variable to solve for?
\answer{The quadratic expression is the expression with the variable term that is raised to the second power; therefore, after completing the square, you solve for the other variable in the equation.}
Q: After completing the square, what properties are helpful in solving the equation?
\answer{the properties of equality}
\end{te-addendum}

% Source page 20: 449, Concept Summary
\begin{te-addendum}{ConceptSummary}{PurposefulQuestions}
\textbf{Pose Purposeful Questions}
Q: What causes the parabola to have its curve on each side of the vertex? Explain in terms of distance.
\answer{The focus and the directrix; any point (P) on a parabola is equidistant from the focus and the directrix. For a vertical parabola, as (x) approaches both negative and positive infinity, the vertical distance increases by the same amount as the horizontal distance, thus creating a curve. For a horizontal parabola, as (y) approaches both negative and positive infinity, the horizontal distance increases by the same amount as the vertical distance, thus creating a curve.}
\te{Printed equal increases in horizontal and vertical distance; likely intended equality of the distance to the focus and the perpendicular distance to the directrix. Coordinate increments are not generally equal.}
\end{te-addendum}
\begin{te-addendum}{ConceptSummary}{CommonError}
\textbf{Common Error: Exercise 7}
Some students may use the incorrect equation (y=\frac{1}{4c}x^2) instead of (x=\frac{1}{4c}y^2). Show them that the focus defines which equation to use. If the coordinates of the focus are of the form ((0,c)), then they should use (y=\frac{1}{4c}x^2) and if the coordinates of the focus are of the form ((c,0)), then they should use (x=\frac{1}{4c}y^2).
\end{te-addendum}
\begin{concept-summary}
\textbf{CONCEPT SUMMARY Parabolas}
DEFINITION A parabola is the set of points on a plane that are equidistant from a given point, the focus, and a given line, the directrix.
GRAPHS The parabolas below have vertex ((0,0)). Parabolas can also be translated anywhere in the coordinate plane.
Vertical: Axis of symmetry: (x=0); Focus: ((0,c)); Directrix: (y=-c).
% FIG 20 768 375 881 488
\placeholder{graph}{Vertical parabola with vertex (0,0) on unlabeled black double-arrow coordinate axes without ticks or grid. Red upward-opening parabola with arrowheads at upper ends; green focus above origin on vertical axis, dashed blue horizontal directrix below origin, vertical axis is axis of symmetry.}
Horizontal: Axis of symmetry: (y=0); Focus: ((c,0)); Directrix: (x=-c).
% FIG 20 955 375 1069 488
\placeholder{graph}{Horizontal parabola with vertex (0,0) on unlabeled black double-arrow coordinate axes without ticks or grid. Red left-opening parabola with arrowheads at left ends; green focus left of origin on horizontal axis, dashed blue vertical directrix right of origin, horizontal axis is axis of symmetry.}
EQUATIONS For a parabola with vertex at the origin and focal length (|c|):
(y=\frac{1}{4c}x^2)
(x=\frac{1}{4c}y^2)
\textbf{Do You UNDERSTAND?}
1. ESSENTIAL QUESTION What are the geometric properties of a parabola, and how do they relate to algebraic representations of a parabola?
\answer{The fixed distance of all points of a parabola is determined by the directrix line and the focus point.}
\te{Printed fixed distance of all points; likely intended that each point is equally distant from the focus and directrix. This common distance varies along the parabola.}
2. Vocabulary Explain what is meant by a conic section.
\answer{A conic section is the curved shape that results from taking a cross-section of a double right cone.}
3. Error Analysis Nicky said that a parabola with the equation (y^2=-9x) has the y-axis as its line of symmetry and opens downward. Explain and correct Nicky's error.
\answer{The parabola (y^2=-9x) has the x-axis as its line of symmetry, and it opens to the left.}
4. Reason What are the focus and directrix of the parabola (y=-4x^2)? Explain how you know.
\answer{The parabola opens downward and the vertex is ((0,0)). Since (a=-4), (c=-\frac{1}{16}), the focus is ((0,-\frac{1}{16})), and the directrix is (y=\frac{1}{16}).}
5. Generalize Explain how the distance from a point to a line is measured.
\answer{It is measured by the perpendicular from the point to the line.}
\textbf{Do You KNOW HOW?}
Write an equation for the parabola with the given focus and directrix.
6. focus: ((0,2)); directrix: (y=-2)
\answer{(y=\frac{1}{8}x^2)}
7. focus: ((-1,0)); directrix: (x=1)
\answer{(x=-\frac{1}{4}y^2)}
8. A parabola has focus ((0,3)) and directrix (y=-3).
a. What is the vertex of the parabola?
\answer{((0,0))}
b. Is the equation in the form (y=ax^2) or (x=ay^2)?
\answer{(y=ax^2)}
c. What is the focal length?
\answer{3}
d. What is the equation of the parabola?
\answer{(y=\frac{1}{12}x^2)}
9. Find the focus and directrix of the parabola that has equation (0=x^2-6x-y+8).
\answer{focus: ((3,-0.75)); directrix: (y=-1.25)}
\end{concept-summary}

% Source page 21: 450, Practice & Problem Solving
\begin{te-addendum}{Practice}{GeneralizeETP}
\textbf{STEP 3 Practice & Problem Solving}
Practice & Problem Solving and video tutorials are available at SavvasRealize.com.
Lesson Practice You may opt to have students complete the automatically scored Practice and Problem Solving items online powered by MathXL for School.
Choose from: Lesson Practice; Adaptive Practice.
You may also take advantage of the bank of exercises for assigning additional practice.
\textbf{Assignment Guide}
\begin{tabular}{ll}
On-Level & Advanced \\
10--12, 16--27 & 10--27 \\
\end{tabular}
\textbf{Engage Through Student Choice}
Promote student agency by allowing students to choose practice items. You may structure this choice in many ways.
For example:
Assign each section a point value. Students choose at least one item from each section and items chosen should have a minimum of 20 total points.
\begin{tabular}{ll}
Understand, Apply & 2 points each \\
Practice & 1 point each \\
Assessment Practice & 1 point each \\
Performance Task & 3 points \\
\end{tabular}
Scan the QR code to access a video tutorial.
\end{te-addendum}
\begin{item-analysis}
\begin{tabular}{lll}
\toprule
Example & Items & DOK \\
\midrule
1 & 16, 17 & 1 \\
1 & 11 & 2 \\
2 & 18 & 1 \\
2 & 13, 26 & 2 \\
3 & 10, 19 & 1 \\
4 & 20 & 2 \\
4 & 23, 24, 27 & 3 \\
5 & 12, 14, 21, 22, 25 & 2 \\
5 & 15 & 3 \\
\bottomrule
\end{tabular}
\end{item-analysis}
\begin{practice}{10}{1}
UNDERSTAND. Use Structure Write an equation of the parabola shown in the graph.
% FIG 21 509 319 638 446
\placeholder{graph}{Unit grid with x and y axes, O at origin, window [-5,5] on both axes, labels -4,-2,2,4, black double-arrow axes. Red upward-opening parabola y=x^2/16 with vertex (0,0) and arrowheads at ends near (-5,1.56) and (5,1.56); red filled focus (0,4) labeled F.}
\answer{(y=\frac{1}{16}x^2)}
\end{practice}
\begin{practice}{11}{2}
Generalize What determines the type, size, and shape of a conic section?
\answer{the way in which the plane intersects the double-right cone}
\end{practice}
\begin{practice}{12}{2}
Error Analysis Describe and correct the error a student made in identifying the vertex and directrix of the parabola with equation (0=2x+y^2-8y+16).
(0=2x+y^2-8y+16)
(-2x=(y-4)^2)
(x=-\frac{1}{2}(y-4)^2)
Vertex: ((4,0)); directrix: (x=-2)
\te{The student's work is marked with a red X.}
\answer{The student wrote a correct equation, but interpreted it incorrectly. The vertex is ((0,4)), not ((4,0)), and since the parabola opens left and (c=-\frac{1}{2}), the directrix is (x=\frac{1}{2}), not (x=-2).}
\end{practice}
\begin{practice}{13}{2}
Look for Relationships What is the shape of a parabola whose focus is very near the directrix?
\answer{The parabola will appear very long and narrow.}
\end{practice}
\begin{practice}{14}{2}
Reason The equation of a parabola is ((x+2)^2=12y). What are the y-intercept(s), if any, of the parabola?
\answer{The y-intercept is ((0,\frac{1}{3})).}
\end{practice}
\begin{practice}{15}{3}
Higher Order Thinking The parabola with equation (y=(x+3)^2-5) has its vertex at ((-3,-5)) and passes through ((1,11)).
a. Identify the focus and directrix of the parabola.
\answer{focus: ((-3,-4.75)); directrix: (y=-5.25)}
b. Name the line of symmetry.
\answer{(x=-3)}
c. Find another point through which the parabola passes that is the same distance from the line of symmetry as ((1,11)).
\answer{((-7,11))}
\end{practice}
\begin{practice}{16}{1}
PRACTICE. Write an equation for all points equidistant from the given point and line. SEE EXAMPLE 1
point ((0,-2)) and line (y=2)
\answer{(y=-\frac{1}{8}x^2)}
\end{practice}
\begin{practice}{17}{1}
Write an equation for all points equidistant from the given point and line. SEE EXAMPLE 1
focus ((\frac{1}{4},0)) and directrix (x=-\frac{1}{4})
\answer{(x=y^2)}
\end{practice}
\begin{practice}{18}{1}
A parabola has a focus of ((-5,0)) and a directrix at (x=5). SEE EXAMPLE 2
a. What is the vertex of the parabola?
\answer{((0,0))}
b. Is the equation in the form (y=ax^2) or (x=ay^2)?
\answer{(x=ay^2)}
c. What is the focal length?
\answer{focal length (=5)}
d. What is the equation of the parabola?
\answer{(x=-\frac{1}{20}y^2)}
\end{practice}
\begin{practice}{19}{1}
What is the equation of the parabola shown in the graph? SEE EXAMPLE 3
% FIG 21 819 547 948 675
\placeholder{graph}{Square grid with spacing 2, window [-10,10] on x and y axes; labels -8,-4,4,8, O at origin, black double-arrow axes x and y. Red downward-opening parabola y=-x^2/32, vertex (0,0), arrowheads at ends near (+/-10,-3.125), red filled focus (0,-8) labeled F.}
\answer{(y=-\frac{1}{32}x^2)}
\end{practice}
\begin{practice}{20}{2}
The cross section of a newer solar cooker is a parabola modeled by the equation (y=\frac{1}{160}x^2), with distances measured in centimeters. The bracket that holds the pan for the cooker needs to be placed at the focus. Assuming that the vertex of the parabolic dish is at the point ((0,0)), at what coordinates should the bracket be placed? SEE EXAMPLE 4
\answer{((0,40))}
\end{practice}
\begin{practice}{21}{2}
Complete the square to find the vertex form, and identify the vertex, focus, and directrix of the parabola with the given equation. SEE EXAMPLE 5
a. (0=-x+2y^2-12y+19)
\answer{vertex ((1,3)); focus ((\frac{9}{8},3)); directrix: (x=\frac{7}{8})}
b. (x^2+24y-8x=-16)
\answer{vertex ((4,0)); focus ((4,-6)); directrix (y=6)}
\te{Printed answer lists the vertex, focus, and directrix, but omits the requested vertex-form equations.}
\end{practice}

% Source page 22: 451, Apply and Assessment Practice
\begin{practice}{22}{2}
APPLY. Make Sense and Persevere The graphs of the equations (0=y^2-8y-12x+28) and (0=y^2-8y+8x+48) have the same directrix.
a. What is the equation of the directrix?
\answer{(x=-2)}
b. What is the distance between the foci?
\answer{10 units}
\end{practice}
\begin{practice}{23}{3}
Model With Mathematics A parabolic mirror has a focus that is located above the vertex of the mirror at the distance shown below. Assume the vertex is at the origin. Write an equation of the parabola that models the cross section of the mirror.
% FIG 22 689 414 816 602
\placeholder{photo}{Parabolic mirror with a marked focus F above its vertex and a vertical bracket labeled 3 in. measuring vertex-to-focus distance.}
\answer{(y=\frac{1}{12}x^2)}
\end{practice}
\begin{practice}{24}{3}
Reason The EuroDish was developed to provide electricity in remote areas, collecting sunlight with a parabolic reflector, which in turn powers the engine located at the focus.
% FIG 22 557 677 769 891
\placeholder{photo}{EuroDish parabolic solar reflector outdoors. Callout Engine points to the receiver at upper left; Vertex: (0,0) points to center of dish at right. White double-arrow segment from vertex to engine is labeled 12 ft.}
a. Write an equation for the parabola.
\answer{Answers may vary. Sample: (y=\frac{1}{48}x^2)}
b. Suppose the diameter of the dish is 24 ft. How deep is the dish?
\answer{3 feet}
\end{practice}
\begin{practice}{25}{2}
ASSESSMENT PRACTICE. Which of the following are true about the graph of the parabolic equation (x+y^2=2y-1)? Select all that apply.
A. opens downward
B. vertex ((1,0))
C. directrix (x=\frac{1}{4})
D. focus ((-\frac{1}{4},1))
\answer{C, D}
\end{practice}
\begin{practice}{26}{2}
SAT/ACT Which graph could have a focus at ((h,0)) and directrix (x=-h), for (h<0)?
A.
% FIG 22 877 414 990 506
\placeholder{graph}{Choice A: black double-arrow x and y axes without ticks or grid, red left-opening parabola whose vertex is right of the origin on the x-axis; arrowheads on both ends at left.}
B.
% FIG 22 1006 414 1131 506
\placeholder{graph}{Choice B: black double-arrow x and y axes without ticks or grid, red right-opening parabola whose vertex is right of the origin on the x-axis; arrowheads on both ends at right.}
C.
% FIG 22 877 520 990 613
\placeholder{graph}{Choice C: black double-arrow x and y axes without ticks or grid, red left-opening parabola with vertex at the origin; arrowheads on both ends at left.}
D.
% FIG 22 1006 520 1131 613
\placeholder{graph}{Choice D: black double-arrow x and y axes without ticks or grid, red right-opening parabola with vertex at the origin; arrowheads on both ends at right.}
\answer{C}
\end{practice}
\begin{practice}{27}{3}
Performance Task The cross section of a flashlight reflector is parabolic. It has an axis of symmetry that is horizontal, and the measure of the widest part of the parabola has the measure shown. The light bulb is at the focus of the parabolic cross section, and the distance from the vertex to the light bulb measures 1 cm.
% FIG 22 865 768 1077 829
\placeholder{illustration}{Cutaway flashlight with teal case, orange batteries at right, gray parabolic reflector opening to left, light bulb at focus on its horizontal axis. Vertical double-arrow dimension at left opening is 6 cm.}
Part A What is the equation of the parabolic cross section of the reflector?
\answer{(x=-\frac{1}{4}y^2)}
Part B How many centimeters deep is the reflector?
\answer{(\frac{9}{4}) cm deep}
\end{practice}

% Source page 23: 451A, Assess & Differentiate
\begin{te-addendum}{Assess}{RtISupport}
\textbf{STEP 4 Assess & Differentiate: LESSON QUIZ}
Use the Lesson Quiz to assess students' understanding of the mathematics in the lesson.
Students can take the Lesson Quiz online or you can download a printable copy from SavvasRealize.com. The Lesson Quiz is also available in the Assessment Sourcebook.
\textbf{Item Analysis}
\begin{tabular}{lll}
Item & DOK & Skills Review and Practice \\
1 & 1 & F77 \\
2 & 1 & F77 \\
3 & 1 & F77 \\
4 & 1 & F77 \\
5 & 2 & F77 \\
\end{tabular}
Use the student scores on the Lesson Quiz to prescribe differentiated assignments.
If students take the Lesson Quiz online, it will be automatically scored and appropriate differentiated practice will be assigned based on student performance.
\begin{tabular}{lll}
Intervention & 0--3 points & Reteach to Build Understanding; Mathematical Literacy and Vocabulary; Lesson Virtual Nerd videos \\
On-Level & 4 points & Enrichment \\
Advanced & 5 points & Enrichment \\
\end{tabular}
\end{te-addendum}
\begin{te-addendum}{Assess}{GeneralizeETP}
\textbf{9-1 Lesson Quiz: Parabolas}
1. What is the equation of a parabola with focus ((0,-8)) and directrix (y=8)?
A. (x=\frac{1}{32}y^2)
B. (x=-\frac{1}{32}y^2)
C. (y=\frac{1}{32}x^2)
D. (y=-\frac{1}{32}x^2)
\answer{1. D}
2. What is the equation of the parabola with vertex ((0,0)) and directrix (x=-\frac{1}{2})?
A. (x=-\frac{1}{2}y^2)
B. (x=\frac{1}{2}y^2)
C. (x=2y^2)
D. (x=-2y^2)
\answer{2. B}
3. What is the equation of the parabola shown in the graph?
% FIG 23 1015 484 1115 585
\placeholder{graph}{Quiz 3: unit square grid, black double-arrow x and y axes, O at origin, window [-5,5] both axes, labels -4,-2,2,4. Black left-opening parabola with vertex (0,0), arrowheads near (-4.2,+/-5), black filled focus F at (-1.5,0). Graph corresponds to x=-y^2/6, although keyed choice C prints a positive coefficient.}
A. (x=-\frac{3}{2}y^2)
B. (x=-\frac{1}{4}y^2)
C. (x=\frac{1}{6}y^2)
D. (x=-\frac{1}{16}y^2)
\answer{3. C}
\te{Printed choice C has positive (\frac{1}{6}); likely (-\frac{1}{6}), consistent with the left-opening graph and focus ((-1.5,0)).}
4. Which statement is true about the parabola with equation (y=\frac{1}{4}(x-3)^2-6)?
A. The focus is ((3,-5)).
B. The vertex is ((-3,-6)).
C. The directrix is (y=-5).
D. It opens to the right.
\answer{4. A}
5. Identify the vertex, focus, and directrix of the parabola with equation (x+y^2+4y-1=0). (Hint: Complete the square to find the vertex form of the equation.)
vertex: (\underline{\hspace{1cm}},\underline{\hspace{1cm}})
focus:
A. ((3\frac{3}{4},-2))
B. ((4\frac{3}{4},-2))
C. ((5\frac{1}{4},-2))
D. ((4\frac{1}{4},-2))
directrix:
A. (x=4\frac{1}{4})
B. (y=4\frac{1}{2})
C. (x=5\frac{1}{4})
D. (y=5\frac{1}{2})
\answer{5. vertex: ((5,-2)); focus: B; directrix: C}
\end{te-addendum}

% Source page 24: 451B, Differentiated Resources
\begin{te-addendum}{Assess}{RtISupport}
\textbf{DIFFERENTIATED RESOURCES: Reteach to Build Understanding}
I = Intervention; O = On-Level; A = Advanced.
Reteach to Build Understanding: I. Provides scaffolded reteaching for the key lesson concepts. MathXL for School.
\textbf{9-1 Reteach to Build Understanding: Parabolas}
A parabola is the set of all points on a plane that are equidistant from a given point, called the focus, and a given line, called the directrix.
The formula for the focal length (c) is (|a|=\frac{1}{4c}). The focal length (|c|) is the distance from the focus to the vertex, measured along the axis of symmetry. If the equation is written in the form (y=ax^2), or (x=ay^2), the vertex is the origin.
\te{Printed (|a|=\frac{1}{4c}); likely (|a|=\frac{1}{4|c|}) when (c) is signed as in the focus coordinates.}
\begin{tabular}{ll}
If the axis of symmetry is the x-axis, then: & If the axis of symmetry is the y-axis, then: \\
(x=\frac{1}{4c}y^2) & (y=\frac{1}{4c}x^2) \\
Focus: ((c,0)) & Focus: ((0,c)) \\
Directrix: (x=-c) & Directrix: (y=-c) \\
\end{tabular}
1. Find the vertex, focus, and directrix of the parabola (y=\frac{1}{16}x^2). Sketch the graph.
Find the focal length (c), which is used to find the focus and directrix, by using (|a|=\frac{1}{4c}), where (\frac{1}{4c}=\frac{1}{16}).
Substitute. (|\frac{1}{16}|=\frac{1}{4c})
Solve for (c). (\frac{1}{16}=\frac{1}{4c})
(4c=16)
(c=\underline{\hspace{1cm}})
Identify the focus. ((0,\underline{\hspace{1cm}}))
Identify the directrix. (y=\underline{\hspace{1cm}})
Identify the vertex. ((\underline{\hspace{1cm}},\underline{\hspace{1cm}}))
\answer{(c=4); focus ((0,4)); directrix (y=-4); vertex ((0,0)).}
% FIG 24 369 537 430 602
\placeholder{graph}{Resource 1: unit square grid with black x and y axes, O at origin, window [-5,5] on both axes, labels +/-2,+/-4. Magenta upward-opening parabola y=x^2/16, focus marked and labeled at (0,4), vertex marked and labeled at (0,0), dashed magenta horizontal directrix y=-4 labeled directrix.}
2. Jordan was given the equation (y=\frac{1}{20}x^2) and found the vertex to be ((0,0)), the focus ((0,-5)) and the directrix to be (y=-5). Find Jordan's mistake and fix it.
\answer{Jordan thought that both the focus and directrix were in the same direction away from the vertex. The focus should be ((0,5)).}
3. Complete the statements given the equation (x=\frac{1}{80}y^2).
a. The blank is ((0,0)).
\answer{vertex}
b. To find the focus use the equation: (|a|=\underline{\hspace{1cm}}).
\answer{(\frac{1}{4c})}
c. The focus is ((\underline{\hspace{1cm}},0)).
\answer{20}
d. The blank is (x=-20).
\answer{directrix}
\end{te-addendum}
\begin{te-addendum}{Assess}{GeneralizeETP}
\textbf{Additional Practice}
I, O. Provides extra practice for each lesson. MathXL for School.
\textbf{9-1 Additional Practice: Parabolas}
Write an equation for a parabola given the focus and directrix.
1. focus ((0,4)) and directrix (y=-4)
\answer{(y=\frac{1}{16}x^2)}
2. focus ((3,0)) and directrix (x=-3)
\answer{(x=\frac{1}{12}y^2)}
3. A parabola has a focus of ((-2,0)) and a directrix at (x=2).
a. What is the vertex of the parabola?
\answer{((0,0))}
b. Is the equation in the form (y=ax^2) or (x=ay^2)?
\answer{(x=ay^2)}
c. What is the focal length?
\answer{2}
d. What is the equation of the parabola?
\answer{(x=-\frac{1}{8}y^2)}
What is the equation of each of the parabolas shown?
4.
% FIG 24 567 521 627 582
\placeholder{graph}{Additional Practice 4: black right-opening parabola with vertex (0,0), focus F at (5,0), arrowheads near (1.25,+/-5). Unit grid, approximate window x [-1,9], y [-5,5], black x/y axes, O at origin, x labels 2,4,6,8, y labels -4,-2,2,4.}
\answer{(x=\frac{1}{20}y^2)}
5.
% FIG 24 693 521 754 582
\placeholder{graph}{Additional Practice 5: black upward-opening parabola y=x^2/28 with vertex (0,0), focus F marked at (0,7), arrowheads near (+/-5,0.9). Unit grid, approximate window x [-5,5], y [-1,9], black x/y axes, O at origin, x labels -4,-2,2,4, y labels 2,4,6,8.}
\answer{(y=\frac{1}{28}x^2)}
6. A TV satellite dish has the shape of a parabola modeled by the equation (x=\frac{1}{120}y^2). The satellite dish has a sensor at its focus. Using the graph of the equation and assuming the base of the dish, or vertex, is positioned at the point ((0,0)), at what coordinates would the sensor be placed?
\answer{((30,0))}
7. Complete the square to identify the vertex, focus, and directrix of the parabola with the equation (0=-y+x^2-6x+2).
\answer{vertex: ((3,-7)); focus: ((3,-6.75)); directrix: (y=-7.25)}
8. What is the value of (c) for the parabola (x=\frac{1}{10}(y+6)^2+2)? Explain.
\answer{The value of (c) is 2.5 because (\frac{1}{4c}=\frac{1}{10}) and (4\times2.5=10).}
9. The midpoint of a pipe with a diameter of 0.5 in. is located 10 in. from a mirror with a parabolic cross section used as a solar collector. The midpoint of the pipe is at the focus of the parabola. Write an equation to model the cross section of the mirror.
\answer{(y=\frac{1}{40}x^2)}
\end{te-addendum}
\begin{te-addendum}{Assess}{RtIExtend}
\textbf{Enrichment}
O, A. Presents engaging problems and activities that extend the lesson concepts. MathXL for School.
\textbf{9-1 Enrichment: Parabolas}
Solve for (a), (c), (h), and (k) in each example. Then add (a+c+h+k). Use the table to match the sum you found in each item with a letter, and then write that letter on the correct blank. Each blank is numbered with the item number.
1. (0=-y-2x^2+32x-135)
\answer{(a=-2), (c=-\frac{1}{8}), (h=8), (k=-7); (a+c+h+k=-1\frac{1}{8}).}
2. (0=-y-8x^2+16x-11)
\answer{(a=-8), (c=-\frac{1}{32}), (h=1), (k=-3); (a+c+h+k=-10\frac{1}{32}).}
3. (0=\frac{1}{8}y^2-\frac{1}{8}y-x-\frac{9}{32})
\answer{(a=\frac{1}{8}), (c=2), (h=-\frac{5}{16}), (k=\frac{1}{2}); (a+c+h+k=2\frac{5}{16}).}
4. (0=3x^2-30x-y-73)
\answer{(a=3), (c=\frac{1}{12}), (h=5), (k=-148); (a+c+h+k=-139\frac{11}{12}).}
5. (0=\frac{x^2}{2}-\frac{x}{2}-y+\frac{1}{2})
\answer{(a=\frac{1}{2}), (c=\frac{1}{2}), (h=\frac{1}{2}), (k=\frac{3}{8}); (a+c+h+k=1\frac{7}{8}).}
6. (0=-\frac{3}{4}y^2+\frac{30}{4}y-x-\frac{103}{4})
\answer{(a=-\frac{3}{4}), (c=-\frac{1}{3}), (h=-7), (k=5); (a+c+h+k=-3\frac{1}{12}).}
7. (0=\frac{5}{4}y^2-\frac{5}{4}y-x-\frac{7}{16})
\answer{(a=\frac{5}{4}), (c=\frac{1}{5}), (h=-\frac{3}{4}), (k=\frac{1}{2}); (a+c+h+k=1\frac{1}{5}).}
8. (0=\frac{3}{8}x^2+\frac{3}{16}x-y+\frac{77}{128})
\answer{(a=\frac{3}{8}), (c=\frac{2}{3}), (h=-\frac{1}{4}), (k=\frac{37}{64}); (a+c+h+k=\frac{263}{192}).}
9. (0=-y-\frac{1}{8}x^2-x+5)
\answer{(a=-\frac{1}{8}), (c=-2), (h=-4), (k=7); (a+c+h+k=\frac{7}{8}).}
\begin{tabular}{lllllllll}
A & B & C & D & E & H & I & K & L \\
(-1\frac{1}{8}) & (-1\frac{1}{5}) & (\uncertain{6\frac{1}{12}}{6\frac{1}{32}}) & (6\frac{11}{12}) & (1\frac{7}{8}) & (-10\frac{1}{32}) & (\uncertain{6\frac{3}{8}}{6\frac{7}{8}}) & (\uncertain{-3\frac{10}{16}}{-3\frac{10}{12}}) & (-1) \\
M & N & O & P & R & S & T & U & Y \\
(2\frac{1}{6}) & (-13\frac{1}{12}) & (\frac{263}{192}) & (-139\frac{11}{12}) & (-3\frac{1}{12}) & (\frac{1}{16}) & (\frac{1}{12}) & (2\frac{7}{8}) & (2\frac{5}{16}) \\
\end{tabular}
What do you call a snake with too much energy?
\answer{A H Y P E R B O A (blanks 1, 2, 3, 4, 5, 6, 7, 8, 9).}
\te{Printed table B is (-1\frac{1}{5}), while item 7 gives positive (1\frac{1}{5}); printed table A is (-1\frac{1}{8}), while item 9 gives (\frac{7}{8}). The completed riddle is retained as printed. Three tiny distractor values remain uncertain.}
\end{te-addendum}
\begin{te-addendum}{Assess}{ELLAddendum}
\textbf{Mathematical Literacy and Vocabulary}
I, O. Helps students develop and reinforce understanding of key terms and concepts.
\textbf{9-1 Mathematical Literacy and Vocabulary: Parabolas}
Match the concept from Column A to the example in Column B. Not all answers will be used.
Column A:
1. Conic section
2. Directrix
3. Focus of a parabola
4. Focal length
5. General form of a 2nd degree equation
6. Parabola
Column B:
(Ax^2+Bxy+Cy^2+Dx+Ey+F=0)
a point such that the distance from this point to any point (P) on a parabola is the same as the distance from (P) to the directrix
a curve formed by the intersection of a plane and a double right cone
(Ax^2+Bx+C=0)
a line such that the distance from this line to any point (P) on a parabola is the same as the distance from (P) to the focus
set of points on a plane that are equidistant from a given point and a given line
the distance between the focus and the vertex
\answer{1. a curve formed by the intersection of a plane and a double right cone; 2. a line such that the distance from this line to any point (P) on a parabola is the same as the distance from (P) to the focus; 3. a point such that the distance from this point to any point (P) on a parabola is the same as the distance from (P) to the directrix; 4. the distance between the focus and the vertex; 5. (Ax^2+Bxy+Cy^2+Dx+Ey+F=0); 6. set of points on a plane that are equidistant from a given point and a given line.}
\end{te-addendum}
\begin{te-addendum}{Assess}{GeneralizeETP}
\textbf{Digital Differentiated Resources}
The Reteach to Build Understanding, Additional Practice, and Enrichment activities are available as digital assignments. These activities are automatically assigned when students complete the lesson quiz online and are automatically scored.
% FIG 24 615 980 1065 1298
\placeholder{illustration}{Tablet showing a digital graphing exercise: Solve the system of equations by graphing, y=3x+4 and y=-2x+4. Use the graphing tool to graph the system. Empty unit grid with x/y axes, window [-10,10] on both axes, labels every 2. Text: Click to enlarge graph; Click the graph, choose a tool in the palette and follow the instructions to create your graph. Question Help menu: Help Me Solve This, View an Example, Video, Textbook, Glossary, Math Tools, Print. Other controls: Exit, 1 part remaining, Clear All, Check Answer, Review progress, Question 5 of 14, Go, Back, Next.}
\end{te-addendum}

\end{document}
